% ============================================================
\chapter{LLM に聞くときの質問の設計}\label{ch:prompt}

BL のビューを LLM から作るとき、$\Omega$ の決め方だけでなく、LLM に「何を、どの形で聞くか」も結果を左右する。この章では、『総合報告書』の実験（D〜G）と先行研究の整理、この資料の結果をつなげて、BL に入れやすい質問の形をまとめる。

\section{本プロジェクトの実験で分かったこと}

『総合報告書』の実験の総合考察は、入力と質問の形を変えた 4 つの実験で、同じ性質がくり返し現れたとまとめている。

\begin{center}
\small
\begin{tabular}{p{0.9cm}p{5.0cm}p{7.0cm}}
\toprule
 & 変えたこと & 予想が沿っていたもの \\
\midrule
既存 & 1 か月分の日次リターンの列を渡す & 渡した月の平均リターン（相関 0.89） \\
D & 日次リターンを要約した指標に置き換え、財務の指標を足す & 直近 1 か月のリターン、12 か月のモメンタム \\
E & 考慮する比率を明言する & 直近 1 か月のリターン。値動きを 0\% にすると成長率 \\
F & 3 ターンの会話にし、答えを 5 段階にする & 直近 1 か月のリターン、または 12 か月のモメンタム（より強く） \\
G & 銘柄の数字を渡さず、リスクの種類ごとに聞く & 市場の直近 1 か月のリターンと VIX。セクターはセクターのモメンタム \\
\bottomrule
\end{tabular}
\end{center}

gpt-4o-mini は、渡された数字の中から「直近の値動き」に当たるものを見つけ、その向きに答える。ほかにも、次のことが分かっている。

\begin{itemize}
\item PER や ROE のような水準の指標は、どの条件でも予想に使われなかった。
\item 30 個の回答はそろいすぎていて、ばらつきは当たり具合を表さない。会話にするとさらにそろい、確率を直接読んでも 1 つの答えに集中する。
\item 銘柄の予想は 76\% がプラスで、リスクプレミアムはほぼ毎月「付く」と答える一方、セクターを指数と比べる問いでは 72\% が「指数より低い」と答えた。問いの形で偏りの向きが変わる。
\end{itemize}

先行研究も同じ方向の結果を示している（『総合報告書』の先行研究の章による）。Chen、Green、Gulen、Zhou（2024）は、LLM の株価の予想が直近の成績に重みを置きすぎ、楽観に偏ることを報告した。数値の時系列を LLM に予想させても良くならないという報告（Tan ほか 2024 など）がある一方、ニュースの見出しやニュースの判定のように、文章を読ませる使い方では成果が報告されている（Lopez-Lira and Tang、Anic ほか 2025）。財務諸表の分析で成果を報告した Kim, Muhn and Nikolaev（2024）は、2025 年 2 月に著者がデータと分析の不整合を理由に取り下げたので、根拠には使わない（2026 年 10 月 19 日に訂正）。

\section{BL の側から見た「良いビュー」}

ここまでの章の結果から、BL に入れやすいビューの条件をまとめる。

\begin{enumerate}
\item \textbf{当たり具合が測れる}：同じ形の予想を、同じ条件で、毎月記録できる。IC と較正の傾きが測れる（\chapref{ch:accuracy}）。
\item \textbf{全体の水準の癖を取り除ける}：楽観・悲観の偏りは全体の水準に出る。相対ビューにできる形がよい（\chapref{ch:shape}）。
\item \textbf{外れの相関が小さい}：銘柄ごとに独立な情報（銘柄ごとのニュースなど）に基づく予想は、外れが相関しにくい。同じ入力（市場の値動き）からすべての銘柄を予想すると、外れが共通になる（\chapref{ch:stacking}）。
\item \textbf{統計モデルにない情報を持つ}：直近の値動きのなぞりは、統計モデルで作れる（そして、この期間はほとんど当たっていない、\chapref{ch:esg}）。LLM を使う価値は、文章のような、統計モデルが扱いにくい情報にある。
\end{enumerate}

\section{聞き方の選択肢と、BL での扱い}

\begin{table}[htbp]
\centering
\caption{LLM への聞き方と、BL での扱い。}
\label{tab:prompt-forms}
\small
\begin{tabular}{p{2.3cm}p{3.2cm}p{3.4cm}p{3.4cm}}
\toprule
聞き方 & 例 & $Q$ の作り方 & $\Omega$ の作り方 \\
\midrule
数値 & 翌月の日次の平均リターン & 平均を引いて相対ビューにし、傾き $b$ で縮める & $\Omega\propto\Sigma$、採用率 $b$ \\
順位・区分 & セクター内の 5 分位、上位・下位 & 区分を点数 $z$ にし、$\IC\cdot\sigma\cdot z$ & 区分の IC から（市場版で $\omega=\sigma^2(1/\IC^2-1)$） \\
確率 & 強気・中立・弱気の確率 & 確率で重み付けた値（\chapref{ch:text}） & エントロピーで順番を付け、大きさは当たり具合で較正 \\
相対の差 & 「この銘柄はセクターを上回るか」 & 差のビュー（$P$ の行の和が 0） & 差の予想の当たり具合から \\
シナリオ & 「関税が上がったら、どのセクターが下がるか」 & 外部の変数へのビュー（\chapref{ch:politics}） & ストレステストとして控えめに（$c\le1$） \\
文章の判定 & 「このニュースはモメンタムの流れを支持するか」 & 統計のビューの採否・大きさの調整 & 判定の当たり具合から \\
\bottomrule
\end{tabular}
\end{table}

\tabref{tab:prompt-forms} のどの聞き方でも、$\Omega$ を LLM 自身に言わせることはしていない。LLM に確信の度合いを言わせると自信過剰になりやすい（Xiong ほか 2024）うえ、本プロジェクトでは回答のばらつきも当たり具合と合わなかった。確信の大きさは、記録した予想と結果から、こちらで測る。

\section{質問を作るときのコツ}

\begin{description}[leftmargin=2em,style=nextline]
\item[直近の値動きだけを渡さない] 渡せば、それをなぞる。値動きのビューが欲しいなら、LLM を通さず統計で作る。LLM には、値動きでは分からない情報（文章）を渡す。
\item[相対で聞く] 「上がるか」ではなく「セクター（や市場）を上回るか」と聞く。全体の水準の楽観・悲観が、答えに入りにくくなる。
\item[数値より順位や区分で聞く] 数値は桁が壊れやすく（実験 E）、単位の取り違えも起きる。区分なら、区分ごとの当たり具合を測って点数にすればよい。
\item[単位と形式をはっきりさせる] 数値で聞くなら、小数か百分率かを明記し、構造化した出力（JSON の形式など）にする。本プロジェクトの元のプロンプトは「percentage change」と書きながら小数を渡していた。
\item[温度と回数を固定し、記録する] 温度、回答の回数、プロンプトの版、モデルの版を、予想と一緒に保存する。どれかを変えたら、当たり具合を測り直す。回答を何回か取って平均するのは、$Q$ の雑音を減らすためで、$\Omega$ を作るためではない。
\item[予想の期間を考える] 1 か月先は、値動きの雑音が大きく、当たり具合が測りにくい。四半期や半年先のほうが、文章の情報が効く余地がある（その分、評価に時間がかかる）。
\item[未来の情報を混ぜない] モデルの知識の期限より後の期間で評価する。「過去の時点のつもりで答えて」という指示では、記憶からの答えを防げないという報告がある（『総合報告書』の先行研究の章）。
\item[複数のモデルは、外れの相関を測ってからまとめる] モデルを増やしても、同じ入力を見ている限り、外れは強く相関する（\chapref{ch:text}）。
\end{description}

\section{記録と評価の仕組み}

BL に LLM のビューを入れるなら、次の量を毎月記録し、一覧で見られるようにしておくとよい。

\begin{center}
\small
\begin{tabular}{p{4.0cm}p{9.6cm}}
\toprule
量 & 何を見るか \\
\midrule
IC（順位相関）と傾き $b$ & 予想が当たっているか。$b$ は採用率の材料 \\
入力との相関 & 予想が入力（直近の値動きなど）のなぞりになっていないか \\
予想の平均と実際の平均 & 全体の水準の癖（楽観・悲観）。相関が負なら全体の水準は入れない \\
回答の分散 ÷ 誤差の 2 乗 & 道具が出す不確かさが当たり具合と合っているか \\
外れの共通部分（$\rho_e$） & 対角の $\Omega$ が許されるか \\
採用率・実質銘柄数・TE & BL がどれだけ動いているか。成績の大部分はここで決まる \\
\bottomrule
\end{tabular}
\end{center}

本プロジェクトの DB には、予想（\code{llm\_views}）、株価、比率（\code{portfolio\_weights}）がそろっている。この表の量は、この資料の計算スクリプトの E、E2、S の節と同じ方法で、毎月計算できる。

\begin{keybox}[この章のまとめ]
\begin{itemize}
\item gpt-4o-mini は、渡された数字の中の直近の値動きをなぞる。予想は楽観に偏り、回答のばらつきは当たり具合を表さない。
\item BL に入れやすいのは、当たり具合が測れ、相対にでき、外れの相関が小さく、統計モデルにない情報を持つビューである。
\item 確信の大きさは LLM に言わせず、記録した予想と結果から測る。
\item 相対で、順位や区分で、単位をはっきりさせて聞き、設定を記録する。
\end{itemize}
\end{keybox}
